

\section{La Guerre des Langues (1902) : L'Humanisme classique contre les "Modernes"}

Au seuil du XXe siècle, la Troisième République française se trouve confrontée à une interrogation
fondamentale qui dépasse largement le cadre strict de la pédagogie pour toucher à l'essence même de
l'identité nationale et de la stratification sociale. La réforme de l'enseignement secondaire de
1902, souvent réduite dans la mémoire collective à une réorganisation technique des cycles d'études,
constitue en réalité l'un des affrontements idéologiques les plus intenses de l'histoire éducative
française. Ce moment historique, cristallisé autour de ce que l'on a coutume d'appeler la « Guerre
des Langues », met en scène une lutte fratricide entre deux visions du monde : l'humanisme
classique, héritier d'une tradition séculaire où le latin et le grec forment l'alpha et l'oméga de
la culture bourgeoise, et les partisans des « Modernes », promoteurs d'un humanisme scientifique et
utilitaire, adapté aux nécessités d'une nation industrielle en pleine mutation.

L'objet de ce rapport est de mener une investigation exhaustive et bibliographique sur ce conflit,
en se concentrant spécifiquement sur la manière dont la réforme de 1902 a redessiné la carte
scolaire et sociale de la France. Il s'agit d'analyser comment, dans un contexte de volonté
démocratique affichée, la création de sections différenciées — et notamment la sanctuarisation de la
Section A (Latin-Grec) — a paradoxalement servi à renforcer les barrières de classe. Loin d'abolir
les privilèges culturels, la réforme a réorganisé la distinction sociale, transformant la maîtrise
des langues anciennes en un capital symbolique d'autant plus précieux qu'il était désormais
optionnel et menacé.

Pour comprendre la portée de cet événement, il convient de replacer la réforme dans la « crise »
multiforme qui frappe l'enseignement secondaire à la fin du XIXe siècle. Cette crise est
diagnostiquée avec une précision chirurgicale par la commission Ribot de 1899, dont les travaux
révèlent l'inadéquation du modèle napoléonien — centralisé, uniforme, et coupé du monde — avec les
aspirations des nouvelles couches sociales.\cite{I-2-1-ref1} La stagnation des effectifs des lycées
publics, l'hémorragie des internes vers l'enseignement privé confessionnel, et la concurrence des
Écoles Primaires Supérieures (EPS) obligent la République à réagir.\cite{I-2-1-ref1} La réponse
politique, portée par le ministre Georges Leygues, sera la tentative d'instaurer une « égalité des
cultures » par la création d'un baccalauréat unique sanctionnant des parcours diversifiés.

Cependant, cette égalité juridique se heurtera violemment à la réalité des hiérarchies sociales. La
Section A, conservatoire des humanités, deviendra le bastion de l'excellence bourgeoise, le lieu où
se reproduit une élite qui refuse de se mêler aux « épiciers » de la section moderne. À travers
l'analyse des débats parlementaires, des polémiques intellectuelles menées par des figures comme
Ferdinand Brunot ou Raoul Frary, et de la sociologie du curriculum développée par André Chervel et
Antoine Prost, nous démontrerons que la Guerre des Langues fut avant tout une guerre de positions
sociales. Nous examinerons comment la rhétorique de la « gymnastique de l'esprit » et de la «
défense de l'âme française » a servi de paravent à une stratégie de distinction, visant à maintenir
une frontière étanche entre ceux qui sont nés pour diriger et ceux qui sont formés pour produire.



\subsection{Le diagnostic de la décadence : l'enquête ribot et la genèse de la réforme}

Avant d'aborder la restructuration de 1902, il est impératif de disséquer le diagnostic posé par la
commission d'enquête parlementaire présidée par Alexandre Ribot en 1898-1899. Cette enquête,
qualifiée de « monumentale » par l'historiographie contemporaine, ne se contente pas de recenser des
dysfonctionnements ; elle dresse le portrait psychologique et sociologique d'une institution en
péril.\cite{I-2-1-ref1}



\subsection{L'enlisement du modèle napoléonien et la crise des effectifs}

À la fin du XIXe siècle, le lycée français vit sur l'acquis du modèle impérial de 1802 : une caserne
monacale dédiée aux humanités classiques, coupée du tissu local, où l'internat joue le rôle de vase
clos pour la formation de l'élite administrative.\cite{I-2-1-ref1} Or, ce modèle s'essouffle. La
commission Ribot met en lumière une stagnation inquiétante des effectifs dans le secondaire public,
alors même que la demande d'éducation s'accroît dans la société.

Le rapport identifie une « hémorragie des internes ».\cite{I-2-1-ref1} Les familles bourgeoises,
traditionnellement clientes des lycées, se détournent de plus en plus de l'internat public, jugé
trop rigide, impersonnel et moralement neutre, au profit des établissements privés ecclésiastiques.
Ces derniers, notamment les collèges jésuites ou marianistes, offrent un encadrement éducatif plus
souple, une attention plus grande à l'individu, et surtout, un environnement socialement homogène et
rassurant.\cite{I-2-1-ref1} La concurrence religieuse n'est pas seulement spirituelle ; elle est
commerciale et sociale. L'enseignement libre a su s'adapter à la demande des familles en proposant
des « maisons d'éducation » là où l'État proposait des « casernes d'instruction ».\cite{I-2-1-ref1}

Parallèlement, sur son flanc gauche, le lycée est attaqué par l'Enseignement Primaire Supérieur
(EPS). Créées par les lois Ferry, ces écoles connaissent un succès foudroyant auprès de la petite
bourgeoisie, des artisans et des employés. Elles offrent un enseignement moderne, sans latin, court
et efficace, débouchant directement sur l'emploi. Le rapport Ribot constate avec amertume que l'EPS
« ne recrute pas dans les rebuts de l'enseignement classique » mais attire des élèves brillants qui
préfèrent l'utile au décoratif.\cite{I-2-1-ref2} Le lycée public se trouve donc pris en étau : trop
« laïc » et « grossier » pour l'aristocratie et la haute bourgeoisie catholique, trop « abstrait »
et « long » pour les classes moyennes ascendantes.



\subsection{Le procès de l'uniformité et de la centralisation}

L'enquête Ribot va plus loin que le simple constat quantitatif. Elle instruit le procès de la
culture administrative française. Les témoins auditionnés, qu'ils soient professeurs, recteurs ou
membres des chambres de commerce, dénoncent l'uniformité étouffante du système. Alexandre Ribot,
républicain libéral, diagnostique un manque de « personnalité » des établissements.\cite{I-2-1-ref1}

Le proviseur, ligoté par la centralisation parisienne, n'est qu'un rouage administratif sans pouvoir
réel sur sa pédagogie ou son recrutement. Les professeurs, fonctionnaires nommés par le Ministère,
ne s'attachent pas à leur établissement ni à leur ville ; ils ne font qu'y passer, l'œil rivé sur
leur avancement et leur mutation vers Paris.\cite{I-2-1-ref1} Cette absence de communauté éducative,
ce manque d'âme des lycées, est perçu comme une faiblesse majeure face aux collèges privés qui
fonctionnent comme des familles.

La réforme de 1902 devra donc répondre à cet impératif : donner de l'autonomie aux lycées pour
qu'ils s'adaptent aux réalités locales et économiques.\cite{I-2-1-ref1} Mais cette décentralisation
administrative ne peut se concevoir sans une refonte des contenus. Si le lycée veut survivre, il
doit s'ouvrir au « moderne », c'est-à-dire aux sciences et aux langues vivantes, et cesser de n'être
que le conservatoire du latin.



\subsection{La parole des acteurs : jaurès, berthelot et la défense de la culture}

Les auditions de la commission Ribot sont un théâtre où s'affrontent les conceptions de la culture.
Jean Jaurès, figure montante du socialisme et normalien classique, y livre une déposition nuancée.
Tout en défendant la nécessité d'ouvrir l'enseignement, il reste attaché à une haute culture
générale et se méfie d'une spécialisation précoce qui enfermerait les enfants du peuple dans un
savoir purement utilitaire.\cite{I-2-1-ref3}

À l'inverse, des scientifiques comme Marcellin Berthelot plaident pour la reconnaissance de la
valeur formatrice des sciences. Ils combattent l'idée reçue selon laquelle seules les lettres
forment l'esprit et les sciences ne fourniraient que des recettes techniques. C'est ici que
s'élabore le concept d'« humanisme scientifique » : la science, par sa rigueur, sa méthode
expérimentale et sa recherche de vérité, est une école de vertu et de logique tout aussi puissante
que la grammaire latine.\cite{I-2-1-ref4}

Le rapport final de la commission Ribot est un compromis politique. Il ne préconise pas la
suppression du latin — ce qui aurait été un suicide politique face à la droite et aux notables —
mais la fin de son monopole. Il propose d'organiser la coexistence pacifique des cultures classique
et moderne au sein d'un même établissement, avec une égalité de dignité. C'est cette proposition qui
servira de socle à la réforme de 1902, transformant le lycée en un système à voies multiples.



\subsection{L'architecture de la réforme de 1902 : une révolution structurelle}

La réforme, portée par le ministre Georges Leygues, est votée par le Parlement en 1902 après des
débats houleux. Elle bouleverse l'organisation séculaire du cursus secondaire en introduisant une
structure en cycles et en sections, préfigurant le lycée moderne.



\subsection{La logique des cycles : différenciation progressive}

La réforme rompt avec le cursus monolithique de la 6ème à la Terminale. Elle divise le secondaire en
deux cycles distincts, chacun ayant une finalité propre.

\begin{itemize}\item \textbf{Le Premier Cycle (de la 6ème à la 3ème) :} Sa fonction est d'assurer les fondements de la culture générale. Cependant, dès l'entrée en 6ème, une bifurcation majeure est opérée. Les élèves sont répartis en deux divisions :\item \textbf{Division A (Classique) :} Le latin est obligatoire dès la 6ème. C'est la filière traditionnelle.\item \textbf{Division B (Moderne) :} Pas de latin. L'enseignement est centré sur le français, les sciences et les langues vivantes renforcées. Cette division est l'héritière de l'enseignement spécial, mais elle est désormais intégrée au cœur du lycée.\cite{I-2-1-ref4}\item \textit{La subtilité du Grec :} En Division A, le grec devient optionnel à partir de la 4ème. C'est une innovation majeure. Jusqu'alors, le bloc \og Latin-Grec\fg{} était indissociable. En rendant le grec facultatif, la réforme crée \textit{de facto} une sous-hiérarchie au sein même du classique : ceux qui font du grec et ceux qui n'en font pas.\cite{I-2-1-ref6}\item \textbf{Le Second Cycle (de la Seconde à la Terminale) :} À l'issue de la 3ème, les élèves ont accès à quatre sections, censées mener à un baccalauréat unique mais comportant des mentions différentes. Cette architecture en quatre colonnes est le cœur du dispositif de 1902.\end{itemize}

\subsection{Les quatre sections : une cartographie de la distinction}

Le Second Cycle structure l'enseignement selon quatre voies, désignées par des lettres (A, B, C, D),
qui vont durablement marquer les esprits et les destins scolaires.\cite{I-2-1-ref4}



\subsection{Tableau 1 : architecture des sections du second cycle (réforme 1902\)}

\begin{table}[h!]\small\centering\begin{tabular}{| l | l | l | l | l |}\hline\textbf{Section} & \textbf{Dénomination} & \textbf{Disciplines Dominantes} & \textbf{Langues Anciennes} & \textbf{Public Cible \& Débouchés (Implicites)} \\ \hline\textbf{Section A} & \textbf{Latin-Grec} & Lettres classiques, Philosophie & Latin (obl.) \+ Grec (obl.) & L'élite traditionnelle, universitaire, juridique, haute bourgeoisie. Voie de l'excellence culturelle. \\ \hline\textbf{Section B} & \textbf{Latin-Langues} & Lettres, Langues Vivantes & Latin (obl.) & Bourgeoisie commerciale, carrières diplomatiques ou administratives. \og Humanités modernes\fg{}. \\ \hline\textbf{Section C} & \textbf{Latin-Sciences} & Sciences Mathématiques \& Physiques & Latin (obl.) & L'élite scientifique (Polytechnique, Centrale), Médecine. L'alliance de la rigueur scientifique et classique. \\ \hline\textbf{Section D} & \textbf{Sciences-Langues} & Sciences, Langues Vivantes & Aucune & Petite bourgeoisie, industrie, commerce. La voie \og moderne\fg{} pure, souvent dévalorisée (\og les épiciers\fg{}). \\ \hline\end{tabular}\end{table}\textit{Analyse des données :} On constate la prédominance structurelle du latin, présent dans trois
sections sur quatre (A, B, C). Seule la section D est totalement affranchie des humanités anciennes.
La Section A se distingue par le cumul des deux langues anciennes, ce qui en fait la section la plus
\og coûteuse\fg{} en temps d'investissement culturel, et donc la plus distinctive
socialement.\cite{I-2-1-ref4}



\subsection{L'égalité des sanctions : le mythe du baccalauréat unique}

La grande promesse de Georges Leygues et des réformateurs est l'« égalité des sanctions ». Le
baccalauréat moderne (issu de la section D) donne désormais les mêmes droits juridiques que le
baccalauréat classique.\cite{I-2-1-ref7} Concrètement, cela signifie qu'un bachelier D peut
s'inscrire en faculté de droit, de médecine ou de lettres, ce qui lui était auparavant interdit (ou
très difficile).

Cette mesure est révolutionnaire sur le papier. Elle vise à casser le monopole de la filière
classique sur les carrières libérales et les fonctions d'État. Elle proclame que la culture
scientifique et moderne a une valeur formatrice équivalente à la culture gréco-
latine.\cite{I-2-1-ref4} Cependant, comme nous le verrons, cette égalité de droit va immédiatement
être contredite par une inégalité de fait, orchestrée par les familles et l'institution elle-même.



\subsection{La guerre idéologique : humanisme classique contre humanisme moderne}

La mise en place de ces structures ne s'est pas faite dans le consensus. Elle a déclenché une
véritable guerre de tranchées intellectuelle, où se sont opposés arguments pédagogiques,
philosophiques et patriotiques.



\subsection{La « question du latin » : raoul frary et les iconoclastes}

La contestation de l'hégémonie latine n'est pas née en 1902. Elle trouve ses racines dans un
pamphlet retentissant publié en 1885 par Raoul Frary : \textit{La Question du
Latin}.\cite{I-2-1-ref9} Frary, journaliste et polémiste, y ose l'impensable : il qualifie le latin
de « préjugé » et de « fétichisme ». Il démonte un à un les arguments traditionnels de la défense
des humanités, affirmant que le temps perdu à traduire péniblement des versions latines (souvent
mal) serait mieux employé à apprendre l'anglais, l'allemand, l'histoire ou les sciences.

Frary est le précurseur des « Modernes ». Pour lui, la France, vaincue en 1870, doit se moderniser
pour survivre face à l'Allemagne industrielle et scientifique. Le maintien artificiel du latin est
vu comme un frein à la puissance nationale. Cette vision utilitariste et nationale est reprise par
les réformateurs de 1902 : l'école doit servir les « besoins économiques du pays » et former des
citoyens adaptés au « milieu moderne ».\cite{I-2-1-ref4}



\subsection{La riposte humaniste : la « gymnastique de l'esprit »}

Face à cette offensive, le camp des « Classiques » — composé de la majorité des professeurs de
lettres, des membres de l'Institut, et de la bourgeoisie conservatrice — affûte ses armes. Ne
pouvant plus défendre le latin sur le terrain de l'utilité pratique, ils se replient sur le terrain
de la psychologie et de la formation mentale.

L'argument central devient celui de la \textbf{« gymnastique de l'esprit »} (ou discipline
formelle). Peu importe que l'on ne parle plus latin, peu importe que l'on oublie ses déclinaisons
après le baccalauréat. Ce qui compte, c'est l'effort fourni. La version latine est présentée comme
l'exercice suprême de logique, d'analyse et de rigueur. Frédéric Queyrat, dans son ouvrage
\textit{La logique chez l'enfant} (1902), se fait l'écho de ces débats, analysant comment l'exercice
de traduction développe les facultés d'abstraction.\cite{I-2-1-ref12}

Pour les humanistes, abandonner le latin, c'est condamner l'esprit français à la mollesse et à
l'approximation. Le latin est l'école de la volonté. Ferdinand Brunot, pourtant favorable à la
réforme, notera ironiquement que cet argument de la « gymnastique » permet de justifier n'importe
quel enseignement inutile, pourvu qu'il soit pénible et difficile.\cite{I-2-1-ref14}



\subsection{La « crise du français » : une panique morale orchestrée}

Vers 1910, la contre-offensive conservatrice prend une nouvelle forme : la dénonciation de la «
Crise du Français ». Des ligues se forment, comme la \textit{Ligue pour la Culture Française}
(présidée par Jean Richepin et soutenue par Poincaré), pour alerter l'opinion : les jeunes Français
ne sauraient plus écrire.\cite{I-2-1-ref15}

Cette baisse de niveau supposée est directement imputée à la réforme de 1902 et à la montée en
puissance des sections sans latin (Section D). On affirme que sans la connaissance des racines
étymologiques latines, la maîtrise de l'orthographe et de la syntaxe française est impossible.
L'argumentaire est habile : il lie la défense du latin à la défense de la langue nationale et de
l'identité française.\cite{I-2-1-ref15}

Ferdinand Brunot, linguiste et historien de la langue, combattra vigoureusement cette idée, montrant
dans ses travaux que le français a sa propre logique et que son enseignement peut être
autonome.\cite{I-2-1-ref17} Pour Brunot, la « crise » est un mythe récurrent agité par les élites
pour bloquer la démocratisation de l'enseignement.



\subsection{La révolution pédagogique et la résistance des corps}

La guerre des langues ne se joue pas seulement sur les programmes, mais aussi sur les méthodes. La
réforme de 1902 tente d'imposer une nouvelle pédagogie pour les langues vivantes, créant une
fracture profonde au sein du corps enseignant.



\subsection{L'imposition de la méthode directe}

Jusqu'en 1902, l'enseignement de l'anglais ou de l'allemand calquait ses méthodes sur celles du
latin : primauté de l'écrit, traduction (thème et version), analyse grammaticale théorique, usage
exclusif du français pour les explications. L'objectif était la culture littéraire, non la
communication.\cite{I-2-1-ref8}

Les instructions de 1902 opèrent une rupture brutale en imposant la \textbf{« Méthode Directe »}.
Celle-ci repose sur plusieurs principes révolutionnaires :

\begin{itemize}\item Interdiction quasi-totale de la langue maternelle en classe.\item Priorité absolue à l'oral et à la prononciation.\item Apprentissage intuitif de la grammaire par l'exemple, sans passer par la traduction.\cite{I-2-1-ref19}\end{itemize}

\subsection{Le malaise des professeurs de langues}

Cette réforme pédagogique, venue d'en haut (du Ministère et des inspecteurs), est vécue comme un
traumatisme par de nombreux professeurs. Formés à l'université classique, agrégés de langues
vivantes sur le modèle des lettres, ils se sentent dépossédés de leur savoir.\cite{I-2-1-ref8}

On leur demande de devenir des « maîtres de conversation », des répétiteurs imitant les méthodes
commerciales (type Berlitz), ce qu'ils jugent indigne de leur statut universitaire. Ils résistent
passivement, continuant à faire de la traduction en cachette ou critiquant l'aspect « superficiel »
des élèves formés à la méthode directe, qu'ils accusent de « jargonner » sans
structure.\cite{I-2-1-ref20}

Cette querelle méthodologique a une conséquence sociologique majeure : elle contribue à dévaloriser
les langues vivantes aux yeux de l'élite scolaire. Puisque la méthode directe est perçue comme
ludique, orale et utilitaire, elle ne peut prétendre à la même dignité intellectuelle que l'austère
version latine. Les langues vivantes deviennent des matières de « service », tandis que le latin
reste la matière de la « culture ».\cite{I-2-1-ref8}



\subsection{André chervel et l'histoire des disciplines}

Les travaux d'André Chervel sur l'histoire des disciplines scolaires éclairent cette période. Il
montre que l'école ne fait pas que transmettre des savoirs savants ; elle produit sa propre culture,
la « culture scolaire ».\cite{I-2-1-ref21} La grammaire scolaire, l'orthographe, l'analyse logique
sont des créations de l'école pour trier les élèves. En 1902, la tentative de remplacer la culture
scolaire du latin (fondée sur l'écrit et la traduction) par une culture scolaire des langues
vivantes (fondée sur l'oral) échoue partiellement parce qu'elle ne fournit pas les mêmes outils de
sélection et de distinction dont le système a besoin.



\subsection{La section a (latin-grec) : bastion sociologique et excellence bourgeoise}

C'est dans l'analyse des flux d'élèves et des représentations sociales que l'on mesure la véritable
issue de la réforme. Malgré l'égalité de droit, la hiérarchie sociale s'est reconstituée
implacablement.



\subsection{La stratégie de la distinction (analyse bourdieusienne)}

Bien que Pierre Bourdieu écrive \textit{La Distinction} et \textit{La Reproduction} bien plus tard,
les mécanismes qu'il décrit sont parfaitement à l'œuvre en 1902.\cite{I-2-1-ref23} Face à la menace
d'une démocratisation (même relative) du secondaire et à l'ouverture des carrières aux bacheliers
modernes, la bourgeoisie met en place une stratégie de repli sur la filière la plus « coûteuse » et
la plus « gratuite » : la Section A.

Le choix du latin et du grec devient un marqueur de classe pur. Apprendre des langues mortes, c'est
signaler que l'on a le temps, que l'on n'est pas pressé par la nécessité de gagner sa vie, que l'on
appartient à une élite héritière. Le grec, devenu optionnel, joue le rôle de « super-distinction » :
c'est le signe de l'excellence absolue, réservé à ceux qui visent l'École Normale Supérieure ou les
plus hautes carrières de l'État.\cite{I-2-1-ref6}



\subsection{La hiérarchie réelle des sections}

Très vite, une hiérarchie informelle mais rigide se met en place dans les lycées :

1. \textbf{Section A (L'Aristocratie) :} Elle recrute les meilleurs élèves, issus des milieux les
plus favorisés. C'est la voie royale. Les professeurs y sont les plus prestigieux.

2. \textbf{Section C (La Haute Bourgeoisie Scientifique) :} Elle attire une élite intellectuelle
forte, celle qui vise Polytechnique. Le maintien du latin y est crucial : il garantit que ces futurs
ingénieurs sont aussi des « hommes cultivés », différents des techniciens.\cite{I-2-1-ref5}

3. \textbf{Section B (L'Entre-deux) :} Souvent perçue comme une voie bâtarde (du latin, mais pas de
grec ; des langues, mais pas assez de sciences).

4. \textbf{Section D (Le Tiers-État Scolaire) :} C'est la section des « Modernes ». Malgré la
qualité des programmes scientifiques, elle subit une dévalorisation symbolique massive. Elle
accueille les élèves jugés inaptes au latin, les fils de commerçants, les boursiers issus du
primaire supérieur.



\subsection{Stéréotypes et stigmatisation : le mythe de l'« épicier »}

La Section D est victime de stéréotypes violents. Ses élèves sont caricaturés en « épiciers » ou en
« cancres ».\cite{I-2-1-ref25} On leur dénie toute finesse d'esprit. Dans l'imaginaire professoral
de l'époque, un élève qui ne fait pas de latin est forcément un esprit borné, utilitariste,
incapable d'abstraction.

Cette stigmatisation a des effets performatifs : les bons élèves, même doués pour les sciences, sont
poussés vers la section C ou A pour ne pas déchoir. La section D devient de facto une filière de
relégation relative, confirmant le préjugé initial. C'est un exemple frappant de prophétie
autoréalisatrice.\cite{I-2-1-ref23}



\subsection{Le paradoxe de la réussite bourgeoise}

Ainsi, la Section A fonctionne comme un « bastion ». Elle protège les enfants de la bourgeoisie de
la concurrence directe avec les enfants des classes moyennes qui affluent vers le moderne. En
imposant le latin comme critère de l'excellence, la bourgeoisie s'assure que la sélection se fait
sur un terrain qu'elle maîtrise culturellement (l'héritage familial, la bibliothèque paternelle) et
non sur le seul terrain des compétences techniques (où les élèves des EPS pourraient rivaliser).



\subsection{L'épilogue : de la réforme bérard à la revanche des sciences}

La tension créée par la réforme de 1902 ne retombe pas. Au contraire, elle s'exacerbe jusqu'à
provoquer une tentative de contre-révolution.



\subsection{La contre-réforme de 1923 : le triomphe éphémère des « ultras »}

En 1923, Léon Bérard, ministre de l'Instruction publique d'un gouvernement de droite (Bloc
National), décide de trancher le nœud gordien. Il supprime les sections modernes au lycée. Par
décret, il rend le latin obligatoire pour tous les élèves de la 6ème à la 3ème.\cite{I-2-1-ref7}

L'argument de Bérard est que l'existence de sections modernes crée des « catégories » de Français et
que seule l'humanité classique peut fonder l'unité nationale. C'est la victoire totale de la ligne
humaniste conservatrice. Pour Bérard, l'égalité républicaine signifie l'accès de tous au latin,
quitte à exclure du lycée ceux qui ne peuvent pas suivre.\cite{I-2-1-ref24}

Cette réforme suscite une levée de boucliers massive de la part de la gauche, des syndicats et des
associations de parents d'élèves modernes. Elle est perçue comme une mesure réactionnaire visant à
fermer le lycée aux classes populaires.



\subsection{Le retour au statu quo et la mutation vers les mathématiques}

Dès la victoire du Cartel des Gauches en 1924, la réforme Bérard est annulée. Les sections modernes
sont rétablies (sous les noms de A' et B, puis Classique et Moderne).\cite{I-2-1-ref7} On revient à
l'architecture de 1902, qui perdurera dans ses grandes lignes jusqu'aux années 1960.

Cependant, une lente mutation s'opère. Au fil du XXe siècle, sous l'effet des besoins technologiques
et de la massification, le critère de l'excellence va se déplacer. La Section C (Latin-Sciences),
puis la future série C (Mathématiques), vont progressivement détrôner la Section A.

Les mathématiques vont remplacer le latin comme outil de sélection et de
distinction.\cite{I-2-1-ref4} Mais il s'agit d'une simple substitution d'instrument : la logique de
la distinction, elle, demeure intacte. La « Section C » des années 1970-1980 jouera exactement le
même rôle sociologique que la Section A de 1902 : celui d'un bastion de l'excellence, protégé par
une barrière d'abstraction infranchissable pour le commun des élèves.

La « Guerre des Langues » de 1902 est un moment fondateur de l'histoire sociale française. Elle ne
peut être réduite à une querelle de programmes. Elle met en lumière la capacité extraordinaire du
système éducatif français — et des élites qui le pilotent ou l'utilisent — à digérer les réformes
démocratiques pour recréer de la distinction.

La réforme de 1902 a tenté d'instaurer un humanisme moderne, égal en dignité à l'humanisme
classique. Elle a échoué sur le plan symbolique. En érigeant la Section A (Latin-Grec) en
conservatoire sacré de la culture « pure » et « désintéressée », la société de la Belle Époque a
signifié que l'utilité sociale et économique (incarnée par les sciences et les langues vivantes de
la section D) restait une valeur subalterne.

Ce conflit a figé pour longtemps les représentations : le littéraire est « cultivé », le
scientifique est « technique », et le moderne est « vulgaire ». Si les contenus ont changé, si le
latin a fini par reculer, la structure mentale héritée de cette guerre — celle d'une hiérarchie
implicite des filières fondée sur leur éloignement de l'utilité immédiate — continue de hanter
l'école française. La Section A de 1902 fut le laboratoire où s'est inventée la méritocratie à la
française : un système ouvert en droit, mais verrouillé de l'intérieur par les codes culturels de la
bourgeoisie.

